\documentclass{article}
\usepackage{amsmath}
\renewcommand{\abstractname}{Synopsis}

\title{Why should I not leave my job to do AI research?}
\author{Cyrus Nolan}
\date{2026-11-22}

\begin{document}

\maketitle

\begin{abstract}

A synopsis

\end{abstract}

My guiding light is to maximize my contribution to humanity. If you're rolling your eyes, I don't care. In my view, the decision variable to maximize over is vocation. Vocation isn't something static that one has no control over. It's something that one can play an active role in shaping and nurturing. The value of a vocation for me is some combination of the following terms:

\begin{itemize}
  \item $C_{\mathrm{ceiling}}(v)$: Contribution ceiling of vocation $v$
  \item $C_{\mathrm{marginal}}(v)$: Marginal contribution of vocation $v$
  \item $C_{\Delta}(v)$: The difference between my marginal contribution and my replacement's in vocation $v$
\end{itemize}

This can be written as an objective function that combines the terms via weighted sum.

\[
  \text{objective}(v) = c_1 \, C_{\mathrm{ceiling}} + c_2 \, C_{\mathrm{marginal}} + c_3 \, C_{\Delta}
\]

My goal can then be restated as shaping vocation to maximize the value function. The weighting of terms is a personal decision. My weights are roughly $c_1 = 0.6$, $c_2 = 0.2$, and $c_3 = 0.2$.

\[
  v^* = \operatorname*{arg\,max}_{v \in \mathcal{V}} \text{objective}(v)
\]

Let's assess the objective function as a function of superintelligence research as the vocation. The ceiling is undeniably high - exactly how high may be up for debate. We can group options into a few buckets.

\end{document}
